Using Theorem 3.3, we can construct the following examples of Lebesgue--Stieltjes measures.

\textbf{Example 3.3.1 (Lebesgue Measure)} \\
Take $F(x) = x$. The resulting Lebesgue--Stieltjes measure defined on $\mathcal{B}(\mathbb{R})$ is known as the Lebesgue measure. We denote this measure by $\lambda$. This measure has the property that $\lambda([a, b]) = b - a$, which can be derived by approaching the closed interval $[a, b]$ by taking intersection of $(a - 1/n, b]$ for $n \ge 1$. Additionally, the Lebesgue measure is translation-invariant.

\textbf{Example 3.3.2 (Uniform Distribution)} \\
Let $F(x)$ be the function defined by
\[
F(x) = \begin{cases} 
0 & \text{for } x < 0 \\ 
x & \text{for } 0 \le x \le 1 \\ 
1 & \text{for } 1 < x. 
\end{cases}
\]
The resulting Lebesgue--Stieltjes measure is a model for uniform distribution between 0 and 1. For example, for $0 < a < b < 1$, the measure of $(a, b)$ is $b - a$.

\textbf{Example 3.3.3 (Continuous Distribution)} \\
Let $f(x)$ be a pdf, i.e., a Riemann integrable function with $f(x) \ge 0$ for all $x$ and $\int_{-\infty}^\infty f(x) \, dx = 1$. The function $F(x) = \int_{-\infty}^x f(t) \, dt$ is a Stieltjes measure function. We denote the corresponding Lebesgue--Stieltjes measure by $P$, which is defined for all Borel sets in $\mathbb{R}$. In particular, if we pick an open set $(a, b)$, for $a < b$, then the probability measure of $(a, b)$ is equal to $\int_a^b f(t) \, dt$.

\textbf{Example 3.3.4 (Point Mass (Dirac Measure))} \\
Consider the Stieltjes measure function given by
\[
F(x) = \begin{cases} 
1 & \text{if } x \ge x_0 \\ 
0 & \text{if } x < x_0, 
\end{cases}
\]
where $x_0$ is a fixed constant. We have a ``point mass'' at $x = x_0$. Let $P$ denote the resulting probability measure. By writing $\{x_0\}$ as the intersection of $(x_0 - n^{-1}, x_0]$ for $n \ge 1$, we can calculate
\[
P(\{x_0\}) = \lim_{n \to \infty} P((x_0 - n^{-1}, x_0]) = \lim_{n \to \infty} F(x_0) - F(x_0 - n^{-1}) = \lim_{n \to \infty} 1 = 1.
\]
The measure $P$ has the property that $P(A) = 1$ if $x_0 \in A$ and $P(A) = 0$ if $x_0 \notin A$.

\textbf{Example 3.3.5 (Discrete Distribution)} \\
For $i = 1, 2, 3, \dots$, let $p_i$ be positive real numbers satisfying $p_i \ge 0$ and $\sum_{i=1}^\infty p_i = 1$. Let $x_i$ be distinct real numbers for $i = 1, 2, 3, \dots$. We want to create a probability space that has probability $p_i$ on the point $x_i$. We suppose that the numbers $x_i$'s are discrete, meaning that we can find a sufficiently small radius $\epsilon$ such that the intervals $(x_i - \epsilon, x_i + \epsilon)$ are mutually disjoint.

We define a Stieltjes measure function $F$ that is flat except at the points $x_i$'s, and there is a vertical jump of size $p_i$ at $x_i$, for $i = 1, 2, 3, \dots$. Let $P$ be the Lebesgue--Stieltjes measure associated with this Stieltjes measure function. Using the same argument as in the last example, we can verify that
\[
P(\{x\}) = \begin{cases} 
p_i & \text{if } x = x_i \text{ for some } i \\ 
0 & \text{otherwise.} 
\end{cases}
\]
The set $A = \{x_1, x_2, x_3, \dots\}$ has $P$-measure $P(A) = 1$. However, as a countable set, the Lebesgue measure (defined in Example 3.3.1) of this set is 0. The distribution of probability is concentrated on the points $x_i$'s.

\textbf{Example 3.3.6 (Cantor Distribution)} \\
In Example 1.2.3, we obtain the Cantor distribution from an infinite sequence of independent random variables. In this example, we take the Cantor function (which is also known as the Devil's staircase) as the Stieltjes measure function and apply the measure extension theorem to obtain the Cantor distribution.

The Cantor function can be defined iteratively, starting with
\[
F_1(x) = \begin{cases} 
3x/2 & \text{if } 0 \le x \le 1/3 \\ 
0.5 & \text{if } 1/3 < x < 2/3 \\ 
0.5 + 0.5(3x - 2) & \text{if } 2/3 \le x \le 1. 
\end{cases}
\]
This is a piece-wise linear function with two pieces with slope $1/6$ in the intervals $[0, 1/3]$ and $[2/3, 1]$. The function $F_1$ is flat in the open interval $I_1 = (1/3, 2/3)$ (See Fig. 3.1).

Given $F_1(x)$, we replace each of the two slant pieces in $F_1(x)$ by a graph with the same shape as $F_1(x)$ and call the resulting function $F_2(x)$. There are three horizontal parts in $F_2(x)$, and they are $I_1$
\[
I_{21} = (1/9, 2/9), \text{ and } I_{22} = (7/9, 8/9).
\]
In general, we obtain $F_{i+1}(x)$ by replacing the $2^i$ non-horizontal parts by a graph that is similar to the graph of $F_1(x)$. After iteration $i$, we have $2^{i-1}$ additional intervals of length $1/3^i$ on which the function $F_i(x)$ has slope zero. The Cantor function $F(x)$ is the limit that we obtain by taking $i \to \infty$.

It can be shown that the functions $F_i(x)$ for $i = 1, 2, 3, \dots$ converge uniformly. Since each $F_i(x)$ is continuous for all $i$, the limit function $F(x)$ is continuous. Therefore, the Lebesgue--Stieltjes measure $P$ induced by the Cantor function has the property $P(\{x\}) = 0$ for any $x$. On the other hand, $F(x)$ has zero slope almost everywhere. Indeed, the total length of the intervals $I_1, I_{21}, I_{22}, \dots$, on which the Cantor function is flat, has Lebesgue measure
% Fig. 3.1 The first three iterations in the construction of the Cantor function
% [Description: Three plots showing the iterative construction of the Cantor function (Devil's Staircase). 
% F1(x) has one flat section in the middle. F2(x) has three flat sections. F3(x) has more steps.]

\[
\sum_{i=1}^\infty 2^{i-1} \frac{1}{3^i} = \frac{1}{2} \sum_{i=1}^\infty (2/3)^i = 1.
\]
As a result, the complement of the union of intervals $I_1 \cup I_{21} \cup I_{22} \cup \dots$ in $[0, 1]$, denoted by $C$, has Lebesgue measure 0. Nonetheless, the probability measure of $C$ is equal to 1. This is the notorious example of singular distribution with continuous distribution function.

In this section, we consider the construction of Lebesgue--Stieltjes measures on the one-dimensional real number line. The idea can be extended to finite-dimensional Euclidean space $\mathbb{R}^d$, where the Stieltjes measure function will be replaced by a function in $d$ variables. For example, in dimension $d = 2$, we can specify a measure by using a function $F(x, y)$ of two variables $x$ and $y$ and define the measure of a rectangle $(a, b] \times (c, d]$ by $F(b, d) - F(a, d) - F(c, b) + F(a, c)$. Readers interested in further details are referred to [1, Section 1.4] and [4, Section 1.1].